\section{Introduction}
\label{sec:introduction}

Geo-distributed LLM serving expands available GPU capacity, but not every
region is equally useful to every request.
Large-scale private WANs expose heterogeneous inter-region paths
~\cite{b4,swan}, while distributed inference systems must additionally
contend with heterogeneous accelerator, serving, and network conditions
~\cite{helix,cassini,skylb,gorgo,solyx}.
Together with heterogeneous request-level SLOs, these differences make
regional capacity request-dependent: a region that is merely attractive for
one request may be the only SLO-feasible choice for another.
A placement therefore determines not only GPU load and cache locality, but
also the WAN path and remote serving conditions experienced by the request.
Interactive requests may impose tight time-to-first-token (TTFT) and
time-between-token (TBT) bounds, whereas longer-running or multi-stage
requests may instead be governed by end-to-end completion or
quality-of-experience objectives
~\cite{distserve,jitserve,quartz,andes,revisit_slo_goodput}.
Consequently, allocating regional capacity well requires reasoning about both
current serving conditions and the SLO-feasible alternatives available to
each request.

Existing distributed routers improve \emph{where} an arriving request should
execute using signals such as serving load, cache locality, WAN conditions,
request-level SLOs, and adaptive routing feedback
~\cite{preble,dualmap,skylb,gorgo,quartz,lodestar}.
However, these routing abstractions largely couple destination selection with
commitment: once a request is assigned using the currently visible state,
later arrivals generally do not influence that placement before execution
begins.
This coupling creates an additional online allocation problem.
A locally attractive placement for the current request can consume
SLO-feasible capacity that a later, more region-constrained request cannot
replace.
The missing question is therefore not only
\emph{where should this request execute?}, but also
\emph{when should an earlier placement become irreversible?}

A two-request example shows why early commitment can waste SLO-feasible
capacity.
Consider two serving regions, $A$ and $B$.
The first request, $r_1$, can satisfy its SLO in either region, whereas a
later request, $r_2$, can succeed only in $A$.
When $r_1$ arrives first, $A$ may appear faster and therefore be selected by a
locally optimal router.
If $r_2$ arrives shortly afterward, however, this decision has consumed its
only feasible capacity.
Placing $r_1$ in $B$ instead would allow both requests to succeed:
\emph{region $A$ is merely faster for $r_1$, but it is essential for $r_2$}.
Thus, optimizing requests independently can reduce aggregate request-level
SLO attainment even when each individual placement appears reasonable when
it is made.

Delaying commitment can reveal the second request, but observation consumes
the same SLO slack needed for WAN transfer, queueing, and inference.
A fixed-delay late-binding policy can therefore observe more demand before
placing a request, yet every millisecond of added deferral also reduces the
time remaining to satisfy that request's SLO.
We call this tension the \emph{information--slack tradeoff}.
Figure~\ref{fig:motivation} illustrates it empirically:
a short fixed deferral can expose useful arrivals and improve SLO attainment,
but the benefit saturates as the deferral consumes an increasing fraction of
request SLO budgets.
This leaves a design opportunity between immediate one-shot placement and
fixed-delay late binding: make arrived demand visible immediately while
preserving only bounded, request-specific placement revisability.

RAVEL's key insight is that deferred commitment does not require deferred
accounting.
Rather than treating placement as one atomic decision, RAVEL separates when a
request begins influencing capacity planning from when its destination becomes
final.
On arrival, RAVEL obtains replica-specific SLO quotes, selects a tentative
serving replica, and immediately attributes the request's prospective demand
to that replica.
Later arrivals therefore observe earlier demand even while the corresponding
placements remain revisable.
While an eligible request remains Router-owned, newly revealed demand may
trigger bounded replanning and transfer both its tentative ownership and
prospective charge to another replica.
The placement becomes final only after successful engine handoff.
By stopping reassignment at this explicit boundary, RAVEL preserves
pre-handoff flexibility without migrating engine-managed KV state.

Realizing this separation safely and efficiently requires three capabilities.
First, RAVEL must determine which replicas remain SLO-feasible under
heterogeneous WAN delays, engine backlog, Prefill and Decode interference, and
request-specific success criteria.
It addresses this challenge with calibrated, replica-specific SLO quotes
derived from live serving state.
Second, tentative work must affect subsequent decisions without becoming
double-counted or stale after reassignment.
RAVEL therefore maintains transferable prospective accounting whose ownership
and replica-dependent charges move atomically with each tentative placement.
Third, revisability must not require repeatedly optimizing over every
outstanding request or delaying work indefinitely.
RAVEL restricts expensive replanning to a bounded Router-owned cohort and ends
all placement revision at the engine-handoff boundary.
Together, these mechanisms retain useful online flexibility while bounding the
latency and state-management costs introduced by the control plane.

We implement RAVEL as an asynchronous routing layer for
vLLM~\cite{vllm} and evaluate it using five serving replicas across three
clusters connected by physical WAN paths.
Our experiments cover conversational traffic, bursty arrivals, and
completion-oriented demand over six offered-load levels.
At $16\times$ load, RAVEL improves request-level SLO attainment over DualMap
by \textbf{8.8}, \textbf{11.2}, and
\textbf{13.4 percentage points} on LMSYS, Burst, and DeepResearch,
respectively.
At the same load, it reduces P95 TTFT by
\textbf{49.1--80.5\%} relative to DualMap.
The online replanning path remains millisecond-scale, with a measured P95
scheduling-pass latency of 2.10\,ms at Burst $12\times$ load.

In summary, this paper makes three contributions:
\begin{itemize}
    \item
    We identify the \emph{request-dependent value of regional capacity} and
    the resulting \emph{information--slack tradeoff}: committing too early
    ignores demand revealed by later arrivals, whereas waiting too long
    consumes the SLO slack needed for successful execution.

    \item
    We introduce an \emph{immediate-accounting, deferred-commitment}
    abstraction and realize it through replica-specific SLO quotes,
    transferable prospective accounting, bounded pre-handoff replanning, and
    an explicit engine-handoff commitment boundary.

    \item
    We implement RAVEL on vLLM and evaluate it on a physical-WAN deployment,
    demonstrating consistent SLO-attainment improvements across
    conversational, bursty, and completion-oriented workloads while retaining
    millisecond-scale online planning overhead.
\end{itemize}




